% Compilen desde esta carpeta o desde la raíz del repositorio.
\makeatletter
\def\input@path{{../../preamble/}{book/preamble/}}
\makeatother
\input{preamble.tex}

\usepackage[spanish]{babel}
\usepackage{lmodern}
\usepackage{titlesec}
\geometry{a4paper,margin=1.6cm}
\linespread{1.10}

\definecolor{HeadingColor}{HTML}{243B53}
\definecolor{InstructionColor}{HTML}{704214}
\definecolor{ConceptColor}{HTML}{205C50}
\definecolor{LinkColor}{HTML}{0047AB}
\hypersetup{allcolors=LinkColor}
\newcommand{\guidelink}[2]{\href{#1}{\textbf{\underline{#2}}}}
\newcommand{\instruction}[1]{\textcolor{InstructionColor}{\textbf{#1}}}
\newcommand{\concept}[1]{\textcolor{ConceptColor}{\textbf{#1}}}
\newcommand{\checkbox}{%
    \begingroup
    \setlength{\fboxrule}{0.4pt}%
    \setlength{\fboxsep}{2pt}%
    \fbox{\rule{0pt}{0.8ex}\rule{0.8ex}{0pt}}%
    \endgroup}

\setlength{\parskip}{4pt}
\setlength{\tabcolsep}{5pt}
\renewcommand{\arraystretch}{1.15}
\setlist{leftmargin=0.65cm,topsep=4pt,partopsep=0pt,parsep=0pt,itemsep=3pt}
\titleformat{\section}{\large\bfseries\color{HeadingColor}}{}{0pt}{}
\titlespacing*{\section}{0pt}{10pt}{3pt}
\newcolumntype{Y}{>{\raggedright\arraybackslash}X}

\hypersetup{pdftitle={C3--C4. Representen y construyan la primera versión},pdfsubject={Emprendimiento e Innovación, Grado 10, Periodo 2}}

\begin{document}

\begin{center}
{\Large\bfseries\color{HeadingColor} C3--C4. Representen y construyan la primera versión\par}
\vspace{3pt}
Emprendimiento e Innovación \enspace|\enspace Grado 10 \enspace|\enspace Periodo 2\\
Guía 2 de 5
\end{center}

\section*{Objetivo y punto de partida}
\concept{C3.} Describe la propuesta mediante bocetos, diagramas, guiones, mapas de servicio o esquemas funcionales.\\
\concept{C4.} Construye un prototipo de baja fidelidad físico, digital, visual, comunicativo o procedimental.

Recuperen D1 y los H de P1, las alternativas y la selección propia fundamentada en la coevaluación de C1--C2 de P2. Identifiquen qué aprendizaje adoptaron y qué supuesto necesita exploración. C3 comunica la propuesta; C4 la convierte en una primera versión que se puede recorrer.

\section*{Documento acumulativo e integridad de la evidencia}
Continúen el \textbf{documento acumulativo en LaTeX} en \guidelink{https://www.overleaf.com/}{Overleaf} o su editor habitual; entreguen fuente y PDF actualizado. Distingan \textbf{P1} y \textbf{P2}. Enlacen evidencias, decisiones, versiones y resultados por códigos o referencias internas, con fuentes y fechas reales.

Distingan \concept{observaciones y comentarios}, \concept{interpretaciones}, \concept{supuestos} e \concept{ideas}; identifiquen simulaciones y límites. \instruction{Usen IA solo para código LaTeX, formato y organización de contenido propio.} No inventen evidencias.

Actualicen la \textbf{bitácora} por sesión: fecha, aportes, evidencias, compromisos, responsables, plazos, estados y acuerdos. Las guías admiten proyectos físicos, digitales, comunicativos, de servicios o procedimentales. Adapten medios y extensión a la evidencia necesaria; los formatos sugeridos son opciones.

\section*{C3 de P2. Representen la propuesta elegida}
\begin{enumerate}[series=pasos]
\item \textbf{Delimiten lo que deben comunicar.} Identifiquen usuario, necesidad, contexto y relación con la propuesta elegida en C2. Expliquen qué busca lograr la persona y cómo respondería la solución; distingan el funcionamiento previsto de lo ya comprobado.
\item \textbf{Hagan comprensibles las partes o pasos.} Representen componentes, secuencia, acciones del usuario y respuestas de la solución. Pueden usar bocetos anotados, diagramas, pantallas, guiones, mapas de servicio u otro medio pertinente. Incluyan conexiones, entradas y resultados previstos cuando sean necesarios para entender el funcionamiento. El formato debe servir a su proyecto.
\item \textbf{Incorporen y expliquen R1.} Den a la representación un código, título y fecha; inclúyanla legible en el documento acumulativo con una explicación que enlace D1, los H pertinentes y la selección de C2. Señalen la función o interacción que construirán primero, por qué es relevante y qué queda fuera de esta versión.
\end{enumerate}

\section*{Entregables de C3}
\begin{itemize}
\item \textbf{Representación R1 anotada.} Incorpórenla legible en C3, con código, título y fecha; debe mostrar usuario, necesidad, funcionamiento y partes o pasos de la propuesta elegida. Conserven su fuente editable si se elaboró fuera de LaTeX.
\item \textbf{Explicación de R1 y alcance de la primera construcción.} Enlacen D1, los H pertinentes y la selección de C2; expliquen qué función o interacción construirán primero, por qué es relevante y qué queda fuera. Esta evidencia describe la propuesta; registren la construcción por separado en C4.
\end{itemize}

\newpage
\section*{C4 de P2. Materialicen la primera versión de baja fidelidad}
\begin{enumerate}[resume=pasos]
\item \textbf{Traduzcan R1 en algo que se pueda usar o recorrer.} Identifiquen qué partes deben existir para mostrar una función o interacción relevante, qué acción hará la persona y qué respuesta podrá percibir. Elijan recursos accesibles al equipo: por ejemplo, un modelo físico, pantallas enlazadas, una pieza comunicativa, una representación de servicio o un procedimiento ejecutable. No necesitan construir toda la solución final.
\item \textbf{Acuerden cómo iniciar la construcción.} Abran PL1 con tareas vinculadas a R1, responsables y roles, recursos y acceso, plazos y dependencias. Registren las condiciones aún no resueltas. Conserven este punto de partida para el seguimiento de C6: el plan ayuda a construir, pero no sustituye el prototipo de C4.
\item \textbf{Construyan V1 y hagan un recorrido interno.} Ejecuten la función o interacción para revisar que se puede mostrar. Identifiquen partes operativas, partes simuladas y partes pendientes; si alguien del equipo simula una respuesta, expliquen cómo. Anoten dificultades de construcción y decisiones. Este recorrido sirve para revisar la versión dentro del equipo; \textbf{no equivale a la prueba con usuarios, pares o docentes de C7}.
\item \textbf{Documenten la versión real.} Registren fecha, código V1 y su relación con R1. Incluyan fotos, capturas, un registro de ejecución, una secuencia anotada u otro soporte que muestre la acción y su respuesta. Si el prototipo queda fuera del PDF, indiquen cómo acceder a él y conserven en el documento evidencia comprensible. Actualicen PL1 y la bitácora con lo realizado y sus pendientes.
\end{enumerate}

\section*{Entregables de C4}
\begin{itemize}
\item \textbf{Prototipo V1 de baja fidelidad disponible.} Entreguen una versión que permita usar o recorrer una función o interacción relevante; si queda fuera del PDF, indiquen cómo acceder a ella.
\item \textbf{Registro de V1 en C4.} Incluyan identificación, fecha, relación con R1 y soportes que muestren la acción y su respuesta, el alcance, las simulaciones y las limitaciones. Debe mostrar lo materializado, no solo una intención o un dibujo sin interacción representable.
\item \textbf{Registro inicial PL1 y enlace a la bitácora.} Muestren tareas vinculadas a R1, responsables y roles, recursos y acceso, plazos, dependencias, trabajo realizado y pendientes; conserven el punto de partida para el seguimiento de C6.
\end{itemize}

\section*{Resultado y continuidad}
Conserven R1 y V1 como referencias de origen. C4 materializa la primera versión; durante la continuación de la construcción, C5 justificará decisiones aplicadas y C6 mostrará cómo organizan y ejecutan el trabajo. Pueden revisar R1 si cambia la propuesta, conservando las versiones y sus motivos.

\section*{Pregunta de reflexión}
¿Qué aprendieron sobre su propuesta al pasar de R1 a V1, a partir de una dificultad o decisión de construcción que puedan respaldar con sus registros?

Incorporen su respuesta al documento acumulativo y susténtenla en la experiencia y las evidencias del equipo.

\end{document}
